%=============================================================================
% DEEP COMMUNICATIONS ANALYSIS: Shannon Capacity, 5D Constellation Theory,
% and Information-Theoretic Security of psi-Field Channels
%
% Extended research paper for Patent #8:
% "psi-Field Communication and Information Systems"
% Inventor: Gary Thomas Alcock
%=============================================================================

\documentclass[twocolumn,superscriptaddress,prd,nofootinbib]{revtex4-2}

%--- Packages ----------------------------------------------------------------
\usepackage{amsmath,amssymb,amsthm}
\usepackage{mathtools}
\usepackage{bm}
\usepackage{physics}
\usepackage{graphicx}
\usepackage{xcolor}
\usepackage{hyperref}
\usepackage{cleveref}
\usepackage{booktabs}
\usepackage{multirow}
\usepackage{array}
\usepackage{siunitx}
\usepackage{tikz}
\usetikzlibrary{arrows.meta,positioning,decorations.pathmorphing,patterns,calc}
\usepackage{pgfplots}
\pgfplotsset{compat=1.18}
\usepackage{tcolorbox}
\tcbuselibrary{theorems,skins}
\usepackage{float}
\usepackage{enumitem}

%--- Custom commands ---------------------------------------------------------
\newcommand{\DFD}{\textsc{DFD}}
\newcommand{\psifield}{\psi}
\newcommand{\astar}{a_*}
\newcommand{\azero}{a_0}
\newcommand{\alphafine}{\alpha}
\newcommand{\ka}{\kappa_\alpha}
\newcommand{\Wfunc}{W}
\newcommand{\mufunction}{\mu}
\newcommand{\Phipot}{\Phi}
\newcommand{\avec}{\mathbf{a}}
\newcommand{\asq}{a^2}
\newcommand{\order}[1]{\mathcal{O}(#1)}
\newcommand{\SNR}{\mathrm{SNR}}
\newcommand{\BER}{\mathrm{BER}}
\newcommand{\GLRT}{\mathrm{GLRT}}
\newcommand{\Prob}{\mathrm{P}}
\newcommand{\Corridor}{\mathcal{C}}
\newcommand{\Template}{\mathcal{T}}
\newcommand{\Fingerprint}{\mathcal{F}}
\newcommand{\Hmin}{H_{\min}}
\newcommand{\Renyientropy}{H_\alpha}
\newcommand{\FisherI}{\mathbf{I}}
\DeclareMathOperator{\Tr}{Tr}
\DeclareMathOperator{\erfc}{erfc}
\DeclareMathOperator{\diag}{diag}
\DeclareMathOperator{\rank}{rank}

\newtheorem{theorem}{Theorem}
\newtheorem{lemma}[theorem]{Lemma}
\newtheorem{proposition}[theorem]{Proposition}
\newtheorem{corollary}[theorem]{Corollary}
\newtheorem{definition}{Definition}

%=============================================================================
\begin{document}

\title{Deep Analysis of $\psi$-Field Communication Channels:\\
Shannon Capacity, 5D Constellation Geometry, and\\
Information-Theoretic Security of Geometry-Locked Authentication}

\author{Gary Thomas Alcock}
\affiliation{Density Field Dynamics Research Programme}

\date{\today}

\begin{abstract}
We present a rigorous information-theoretic analysis of communication
channels based on the scalar refractive field $\psi$ of Density Field
Dynamics. Starting from the complete signal model with quantum vacuum,
thermal, and systematic noise contributions, we derive the Shannon
capacity of a single $\psi$-corridor in closed form, obtaining
$C_\psi = B\log_2(1 + P_\psi/N_0 B)$ with explicit expressions for the
noise floor $N_0$ in terms of fundamental constants. We prove that the
5D composite constellation $(\delta\psi, \dot\psi, \nabla\psi)$ achieves
a capacity gain of $\Gamma_5 = 5\log_2(1 + \SNR/5)/\log_2(1+\SNR)$
over the scalar channel, approaching the factor-of-5 limit at high SNR,
and derive the optimal sphere-packing density in the 5D symbol space.
The central result is a complete information-theoretic security proof for
geometry-locked authentication: we show that inverting the corridor
fingerprint to determine the source geometry is computationally equivalent
to inverting an elliptic PDE --- a problem in complexity class
$\mathrm{NP}$-hard for discrete source configurations. We compute the
min-entropy of the geometric fingerprint as
$\Hmin \geq N_s \log_2(L/\ell_{\rm corr})$ where $N_s$ is the number of
sampling points and $\ell_{\rm corr}$ is the field correlation length,
and prove that an adversary's spoofing probability decays as
$P_{\rm spoof} \leq \exp(-\beta d^2/\ell_{\rm corr}^2)$ where $d$ is
the adversary's displacement and $\beta$ depends on corridor geometry.
Extensive comparison with 15 existing technologies --- including
GPS anti-jam systems, quantum key distribution, neutrino communications,
underwater acoustics, and through-wall radar --- demonstrates that
$\psi$-corridors occupy a unique capability niche. A combined
communications-navigation demonstrator sharing infrastructure is designed,
with explicit resource allocation trade-offs computed via convex
optimization.
\end{abstract}

\maketitle
\tableofcontents

%=============================================================================
\section{Introduction}\label{sec:intro}
%=============================================================================

The companion paper~\cite{Alcock2025PsiComms} established the
foundational framework for $\psi$-field communications: corridor
geometry, composite modulation, basic capacity bounds, and security
concepts. This paper goes substantially deeper on three fronts:
(i)~rigorous derivation of channel capacity from first principles with
complete noise modeling, (ii)~proof that 5D constellation geometry
yields quantifiable gains with explicit sphere-packing analysis, and
(iii)~a full information-theoretic security proof reducing
geometry-locked authentication to the hardness of elliptic PDE inversion.

We also conduct an exhaustive comparison with 15 competing communication
and security technologies, identifying precise capability gaps that
$\psi$-corridors fill, and design a unified communications-plus-navigation
demonstrator.

\subsection{DFD Signal Model Recap}

In Density Field Dynamics~\cite{Alcock2025DFD}, the scalar refractive
field $\psi(\mathbf{r},t)$ satisfies the field equation
\begin{equation}
\nabla \cdot \bigl[\mu(|\nabla\psi|/\astar)\nabla\psi\bigr]
= -\frac{8\pi G}{c^2}(\rho - \bar\rho),
\label{eq:field-eq}
\end{equation}
with optical metric
$d\tilde{s}^2 = -c^2 dt^2/n^2 + d\mathbf{r}^2$, $n = e^\psi$.
For communication applications in the high-gradient regime
($|\nabla\psi|/\astar \gg 1$, $\mu \to 1$), a $\psi$-corridor carries
information encoded in the perturbation
$\delta\psi(\mathbf{r},t)$ about a background $\psi_0$.

%=============================================================================
\section{Complete Noise Model from First Principles}\label{sec:noise}
%=============================================================================

\subsection{Quantum Vacuum $\psi$-Fluctuations}

The $\psi$-field, being a dynamical scalar degree of freedom, exhibits
vacuum fluctuations. From the linearized action for $\psi$-perturbations
around flat background~\cite{Alcock2025DFD}:
\begin{equation}
S_{\delta\psi} = \frac{1}{8\pi G}\int d^4x\left[
\frac{1}{c^2}(\partial_t \delta\psi)^2 - (\nabla\delta\psi)^2
\right],
\label{eq:linearized-action}
\end{equation}
we identify the effective mass parameter
$m_{\rm eff}^2 = 1/(8\pi G)$ (in units where $\hbar = c = 1$).
Restoring dimensions, the vacuum spectral density of $\psi$-fluctuations
in a detector of effective volume $V_d$ is
\begin{equation}
S_\psi^{\rm qm}(f) = \frac{4\hbar G}{c^3 V_d}
\cdot \frac{1}{1 + (f/f_c)^2},
\label{eq:quantum-noise-spectrum}
\end{equation}
where $f_c = c/(2\pi \ell_d)$ is the cutoff frequency set by the
detector size $\ell_d = V_d^{1/3}$.

\begin{proof}[Derivation]
Quantize $\delta\psi$ in the box of volume $V_d$:
\begin{equation}
\delta\hat\psi(\mathbf{r},t) = \sum_\mathbf{k}
\sqrt{\frac{4\pi G \hbar}{c\,\omega_k\,V_d}}
\left(\hat{a}_\mathbf{k} e^{i(\mathbf{k}\cdot\mathbf{r} - \omega_k t)}
+ \text{h.c.}\right),
\end{equation}
where $\omega_k = c|\mathbf{k}|$ and the normalization follows from
$[S_{\delta\psi}] = \hbar$. The two-point function gives
\begin{equation}
\langle 0|\delta\hat\psi(\mathbf{r},t)\delta\hat\psi(\mathbf{r},t)|0\rangle
= \frac{4\pi G\hbar}{c}\int \frac{d^3k}{(2\pi)^3}\frac{1}{\omega_k V_d}.
\end{equation}
Converting to the one-sided power spectral density
$S_\psi^{\rm qm}(f) = 2\int_{-\infty}^\infty R_\psi(\tau)e^{-2\pi i f\tau}d\tau$
and imposing the detector volume as an infrared cutoff yields
Eq.~\eqref{eq:quantum-noise-spectrum}.
\end{proof}

\paragraph{Numerical evaluation.}
For a detector with $V_d = (1\;\text{m})^3$:
\begin{equation}
\sqrt{S_\psi^{\rm qm}} = \sqrt{\frac{4\hbar G}{c^3}}
\approx \sqrt{\frac{4 \times 1.05\times 10^{-34}\times 6.67\times 10^{-11}}
{2.7\times 10^{25}}}
\approx 1.0\times 10^{-35}\;\text{Hz}^{-1/2}.
\end{equation}
This is the fundamental quantum noise floor for $\psi$-field detection.

\subsection{Thermal Noise}

Thermal fluctuations in the detector mass-energy produce stochastic
$\psi$ perturbations via the source term in Eq.~\eqref{eq:field-eq}.
For a detector at temperature $T_d$ with mechanical quality factor $Q_m$
and effective mass $M_d$ coupling to the $\psi$-mode:
\begin{equation}
S_\psi^{\rm th}(f) = \frac{4k_B T_d}{M_d c^4}
\cdot \frac{f_0/Q_m}{(f^2 - f_0^2)^2 + (f f_0/Q_m)^2}
\cdot \frac{G M_d}{c^2},
\label{eq:thermal-noise-full}
\end{equation}
where $f_0$ is the mechanical resonance frequency. On resonance:
\begin{equation}
S_\psi^{\rm th}(f_0) = \frac{4 k_B T_d G Q_m}{c^6 f_0}.
\label{eq:thermal-on-resonance}
\end{equation}
For $T_d = 300\;\text{K}$, $Q_m = 10^6$, $f_0 = 1\;\text{kHz}$:
\begin{equation}
\sqrt{S_\psi^{\rm th}(f_0)} \approx
\sqrt{\frac{4\times 4.1\times 10^{-21}\times 6.67\times 10^{-11}\times 10^6}
{7.3\times 10^{51}\times 10^3}}
\approx 3.5\times 10^{-39}\;\text{Hz}^{-1/2}.
\end{equation}
Thermal noise on resonance is subdominant to quantum vacuum noise by
$\sim 4$ orders of magnitude.

\subsection{Systematic Noise: Gravitational Confusion Background}

The astrophysical gravitational-wave background and local seismic/tidal
signals create a confusion floor. From the $\psi$-GW correspondence
$\delta\psi \sim (c^2/2)\int h\,dt$ (frequency domain:
$\tilde{\delta\psi}(f) \sim c^2 \tilde{h}(f)/(4\pi i f)$):
\begin{equation}
S_\psi^{\rm grav}(f) = \frac{c^4}{16\pi^2 f^2}\,S_h(f),
\label{eq:grav-confusion}
\end{equation}
where $S_h(f)$ is the gravitational-wave strain spectral density.
For the stochastic GW background at $f \sim 1\;\text{Hz}$,
$\sqrt{S_h} \sim 10^{-24}\;\text{Hz}^{-1/2}$
(LISA sensitivity target), giving
\begin{equation}
\sqrt{S_\psi^{\rm grav}(1\;\text{Hz})} \sim
\frac{9\times 10^{16}}{4\pi}\times 10^{-24}
\sim 7\times 10^{-9}\;\text{Hz}^{-1/2}.
\end{equation}
However, seismic isolation and correlation techniques can suppress this
by $>60\;\text{dB}$, and the confusion background is concentrated at
low frequencies ($<10\;\text{Hz}$).

\subsection{Total Noise Spectral Density}

The total noise PSD is
\begin{equation}
N_0(f) \equiv S_\psi^{\rm tot}(f) = S_\psi^{\rm qm}(f) + S_\psi^{\rm th}(f)
+ S_\psi^{\rm grav}(f) + S_\psi^{\rm sys}(f),
\label{eq:N0-total}
\end{equation}
where $S_\psi^{\rm sys}$ captures instrumental systematics (readout
noise, coupling drift, etc.). In the band $10\;\text{Hz} < f < 10\;\text{kHz}$
with seismic isolation, the dominant contribution is quantum vacuum noise:
\begin{equation}
N_0 \approx S_\psi^{\rm qm} = \frac{4\hbar G}{c^3 V_d}
\approx 1.0\times 10^{-69}\;\text{Hz}^{-1}\quad\text{(for }V_d = 1\;\text{m}^3\text{)}.
\end{equation}

%=============================================================================
\section{Complete Shannon Capacity Derivation}\label{sec:capacity}
%=============================================================================

\subsection{Signal Model}

The received $\psi$-field perturbation at the output of a corridor
$\Corridor$ of length $L$ is
\begin{equation}
\delta\psi_{\rm rx}(t) = h_\Corridor \cdot \delta\psi_{\rm tx}(t - \tau_\Corridor)
+ w(t),
\label{eq:signal-model}
\end{equation}
where:
\begin{itemize}
\item $h_\Corridor = e^{-\alpha_\psi L}$ is the corridor attenuation
factor, with $\alpha_\psi$ the $\psi$-field absorption coefficient
(vanishing for the linearized field equation: $\alpha_\psi = 0$),
\item $\tau_\Corridor = (e^{\psi_0}/c)L$ is the propagation delay,
\item $w(t)$ is zero-mean Gaussian noise with PSD $N_0(f)$.
\end{itemize}

A critical feature of $\psi$-corridors: the linearized field equation
is a Laplace/Poisson equation (elliptic), not a wave equation. The
``propagation'' of $\psi$-perturbations is instantaneous in the
quasi-static regime --- there is no dispersive spreading. This means:
\begin{equation}
h_\Corridor = 1\quad\text{(lossless corridor in quasi-static regime)}.
\label{eq:lossless}
\end{equation}
The corridor preserves signal amplitude without distance-dependent
attenuation, fundamentally unlike electromagnetic propagation where
free-space path loss scales as $(4\pi d/\lambda)^2$.

\subsection{Signal-to-Noise Ratio}

Define the transmit ``power'' as the mean-square $\psi$-perturbation rate:
\begin{equation}
P_\psi \equiv \lim_{T\to\infty}\frac{1}{T}\int_0^T
|\delta\psi_{\rm tx}(t)|^2\,dt \qquad [\text{s}^{-1}],
\end{equation}
where $\psi$ is dimensionless so $P_\psi$ has dimensions of inverse time
(equivalently, $\psi^2\cdot$Hz). The SNR per unit bandwidth is
\begin{equation}
\frac{\SNR}{B} = \frac{P_\psi \cdot |h_\Corridor|^2}{N_0 \cdot B}
= \frac{P_\psi}{N_0\,B}.
\label{eq:SNR-def}
\end{equation}

For a corridor-modulated signal with $\psi$-pump power generating
$\delta\psi_{\rm rms} \sim 10^{-20}$ at the receiver and bandwidth
$B = 1\;\text{kHz}$:
\begin{equation}
\SNR = \frac{(\delta\psi_{\rm rms})^2 B}{N_0 B}
= \frac{(10^{-20})^2}{10^{-69}} = 10^{29}.
\label{eq:SNR-numerical}
\end{equation}
This extraordinary SNR reflects the extreme weakness of $\psi$-noise
--- the fundamental quantum limit is set by the Planck-scale
gravitational vacuum fluctuations. The practical SNR will be limited
by the achievable $\psi$-pump amplitude rather than the noise floor.

\subsection{Shannon Capacity: Scalar Channel}

\begin{theorem}[Scalar $\psi$-Channel Capacity]\label{thm:scalar-capacity}
The capacity of a single $\psi$-corridor with bandwidth $B$, transmit
power $P_\psi$, and white Gaussian noise of one-sided spectral density
$N_0$ is
\begin{equation}
\boxed{
C_1 = B\log_2\!\left(1 + \frac{P_\psi}{N_0\,B}\right)
\;\text{bits/s}
}
\label{eq:C1}
\end{equation}
achieved by Gaussian codebook with power spectral density $P_\psi/B$
uniformly distributed across the band.
\end{theorem}

\begin{proof}
The signal model~\eqref{eq:signal-model} with $h_\Corridor = 1$ and
additive white Gaussian noise $w(t)$ with PSD $N_0/2$ (two-sided) is
isomorphic to the scalar AWGN channel. We verify the conditions:
\begin{enumerate}
\item \emph{Linearity:} The corridor is lossless and non-dispersive in
the quasi-static regime, so the channel transfer function is
$H(f) = 1$ for $|f| \leq B/2$.
\item \emph{Gaussianity:} The quantum vacuum noise $w(t)$ is a
zero-mean Gaussian process by the central limit theorem applied to
the superposition of independent vacuum mode contributions.
The thermal and systematic contributions are also Gaussian in the
central-limit regime.
\item \emph{Stationarity:} The noise PSD is time-independent for a
stationary detector.
\item \emph{Power constraint:}
$\mathbb{E}[|\delta\psi_{\rm tx}|^2] \leq P_\psi/B$ per frequency bin.
\end{enumerate}
Shannon's channel coding theorem~\cite{Shannon1948,CoverThomas2006}
then gives the capacity as~\eqref{eq:C1}. The maximum is achieved by
i.i.d.\ Gaussian input symbols at rate $1/T_s$ with $T_s = 1/B$.

For colored noise $N_0(f)$, the water-filling generalization gives
\begin{equation}
C_1 = \int_0^B \log_2\!\left(1 + \frac{P_\psi(f)}{N_0(f)}\right)df,
\end{equation}
with $P_\psi(f) = \max(\nu - N_0(f), 0)$ and $\int P_\psi(f)\,df = P_\psi$.
\end{proof}

\paragraph{Spectral efficiency.}
The spectral efficiency $\eta = C_1/B$ is
\begin{equation}
\eta = \log_2(1 + \SNR)\;\text{bits/s/Hz}.
\end{equation}
For $\SNR = 10^{29}$:
$\eta \approx 29\log_2 10 \approx 96\;\text{bits/s/Hz}$.
This represents an enormous theoretical spectral efficiency, limited
in practice by the achievable $\psi$-pump power and detector dynamic range.

\subsection{Shannon Capacity: 5D Composite Channel}

\begin{theorem}[5D Composite Channel Capacity]\label{thm:5D-capacity}
The capacity of the composite $\psi$-channel with five independent
degrees of freedom $\mathbf{s} = (\delta\psi, \dot\psi, \partial_x\psi,
\partial_y\psi, \partial_z\psi) \in \mathbb{R}^5$, subject to total power
constraint $\sum_{j=1}^5 P_j \leq P_{\rm tot}$ and independent Gaussian
noise with per-dimension PSD $N_j$, is
\begin{equation}
\boxed{
C_5 = \sum_{j=1}^{5} B_j\log_2\!\left(1 + \frac{P_j^*}{N_j\,B_j}\right)
}
\label{eq:C5}
\end{equation}
where $P_j^* = \max(\nu - N_j B_j, 0)$ is the water-filling allocation
with Lagrange multiplier $\nu$ set by $\sum_j P_j^* = P_{\rm tot}$.
\end{theorem}

\begin{proof}
The five observables $(\delta\psi, \dot\psi, \nabla_x\psi, \nabla_y\psi,
\nabla_z\psi)$ constitute independent measurements at each spacetime
point because:
\begin{enumerate}
\item $\delta\psi$ and $\dot\psi$ are the field value and its time
derivative --- canonically conjugate observables in the $\psi$-Hamiltonian.
In the classical limit (high occupation number), they behave as
independent Gaussian channels.
\item $\nabla_i\psi$ for $i=x,y,z$ are spatial derivatives measured by
separated detectors (gradiometers), with noise independent across axes
for isotropic detector noise.
\item Cross-correlations between $\delta\psi$ and $\nabla_i\psi$ vanish
at a single point by parity: $\langle\delta\psi\,\partial_i\delta\psi\rangle
= \frac{1}{2}\partial_i\langle(\delta\psi)^2\rangle = 0$ for
translation-invariant vacuum.
\end{enumerate}
The 5D channel therefore decomposes into five parallel Gaussian channels,
and the total capacity is the sum of individual capacities, optimized by
water-filling~\cite{CoverThomas2006}.
\end{proof}

\subsection{5D Capacity Gain: Rigorous Derivation}

\begin{theorem}[5D Constellation Capacity Gain]\label{thm:gain}
Define the capacity gain ratio
\begin{equation}
\Gamma_5 \equiv \frac{C_5}{C_1}.
\end{equation}
For equal noise levels $N_j = N_0$ and equal bandwidths $B_j = B$:
\begin{equation}
\Gamma_5 = \frac{5\log_2(1 + \SNR_{\rm tot}/(5B N_0))}
{\log_2(1 + \SNR_{\rm tot}/(B N_0))}.
\label{eq:gain-ratio}
\end{equation}
In the high-SNR limit ($\SNR_{\rm tot}/(BN_0) \gg 5$):
\begin{equation}
\Gamma_5 \to \frac{5(\log_2 \SNR_{\rm tot} - \log_2(5BN_0))}
{\log_2 \SNR_{\rm tot} - \log_2(BN_0)}
\to 5 - \frac{5\log_2 5}{\log_2(\SNR_{\rm tot}/(BN_0))},
\label{eq:gain-highSNR}
\end{equation}
approaching 5 asymptotically. In the low-SNR limit ($\SNR \ll 1$):
\begin{equation}
\Gamma_5 \to \frac{5\cdot\SNR/(5BN_0)}{\ \SNR/(BN_0)\ } = 1,
\end{equation}
since all channels are power-starved and the gain vanishes.
\end{theorem}

\begin{proof}
With equal allocation $P_j = P_{\rm tot}/5$, each sub-channel has
$\SNR_j = P_{\rm tot}/(5BN_0)$. Thus
$C_5 = 5B\log_2(1 + \SNR/5)$ where $\SNR = P_{\rm tot}/(BN_0)$, and
$C_1 = B\log_2(1 + \SNR)$. The ratio is~\eqref{eq:gain-ratio}.
Taylor expanding $\log_2(1+x) = x/\ln 2 + \order{x^2}$ for small $x$
gives $\Gamma_5 \to 1$. For large $x$,
$\log_2(1+x) \approx \log_2 x$, giving
$\Gamma_5 \approx 5\log_2(\SNR/5)/\log_2(\SNR)
= 5(1 - \log_2 5/\log_2\SNR) \to 5$.
\end{proof}

\paragraph{Numerical example.}
For $\SNR = 10^{10}$ (feasible with amplified $\psi$-emitters):
\begin{equation}
\Gamma_5 = \frac{5\log_2(2\times 10^9)}{\log_2(10^{10})}
= \frac{5\times 30.9}{33.2} = 4.65.
\end{equation}
The 5D channel achieves 93\% of the theoretical $5\times$ limit.

\subsection{Sphere-Packing in the 5D Symbol Space}

The constellation gain from higher-dimensional modulation is understood
through sphere-packing theory. The maximum number of non-overlapping
$\epsilon$-balls in a region of volume $V$ in $\mathbb{R}^d$ scales as
\begin{equation}
M_{\max} \sim \frac{V}{V_d(\epsilon)} = \frac{V}{\pi^{d/2}\epsilon^d/\Gamma(d/2+1)}.
\end{equation}

\begin{proposition}[5D Packing Advantage]
For a total energy constraint $E_{\rm tot} = M\bar{E}_s$ where
$\bar{E}_s$ is the average symbol energy, and minimum distance
$d_{\min}$ for a target BER, the number of distinguishable symbols in
$d$ dimensions is
\begin{equation}
\log_2 M_d \leq \frac{d}{2}\log_2\!\left(\frac{2e\bar{E}_s}{d\,d_{\min}^2}\right)
+ \frac{d}{2}\log_2\!\left(\frac{d}{2\pi e}\right) + o(d).
\label{eq:sphere-packing}
\end{equation}
The gain from $d=5$ versus $d=1$ at equal energy and BER is
\begin{equation}
G_{\rm pack}^{(5/1)} = \frac{\log_2 M_5}{\log_2 M_1}
\approx \frac{5}{2}\cdot\frac{\log_2(2e\bar{E}_s/(5d_{\min}^2))
+ \log_2(5/(2\pi e))}
{\frac{1}{2}\log_2(2e\bar{E}_s/d_{\min}^2)}.
\label{eq:packing-gain}
\end{equation}
For $\bar{E}_s/d_{\min}^2 = 10^4$ (high SNR):
$G_{\rm pack}^{(5/1)} \approx 4.3$.
\end{proposition}

\begin{proof}
Apply the Minkowski-Hlawka theorem for sphere packing in
$\mathbb{R}^d$: the densest lattice packing satisfies
$\delta_d \geq \zeta(d)/V_d(1)$ where $\zeta$ is the Riemann zeta
function and $V_d(1)$ is the unit ball volume. The maximum number of
points at minimum distance $d_{\min}$ inside a $d$-ball of radius
$R = \sqrt{d\bar{E}_s}$ (energy shell for uniform constellation) is
$M_d \leq (R/d_{\min}+1)^d \cdot \delta_d$. Taking logarithms and
using Stirling's approximation yields~\eqref{eq:sphere-packing}.
\end{proof}

%=============================================================================
\section{Geometry-Locked Authentication: Full Security Proof}
\label{sec:security}
%=============================================================================

\subsection{Formal Security Model}

\begin{definition}[Corridor Fingerprint Space]
The corridor fingerprint space $\mathcal{F}_\Corridor$ is the
Hilbert space $L^2([0,L]; \mathbb{R}^5)$ of square-integrable
5-component vector functions along the corridor path $[0,L]$:
\begin{equation}
\Fingerprint_\Corridor(s) = \bigl(\delta\psi(s),\;\dot\psi(s),\;
\partial_x\psi(s),\;\partial_y\psi(s),\;\partial_z\psi(s)\bigr).
\end{equation}
The inner product is
$\langle F_1, F_2\rangle = \int_0^L F_1(s)\cdot F_2(s)\,ds$.
\end{definition}

\begin{definition}[Source-to-Fingerprint Map]
The source-to-fingerprint map $\Phi: \mathcal{D} \to \mathcal{F}_\Corridor$
takes a source density distribution $\delta\rho \in \mathcal{D}
\subset L^2(\mathbb{R}^3)$ to the corridor fingerprint:
\begin{equation}
\Phi[\delta\rho](s) = \Fingerprint_\Corridor[\delta\psi[\delta\rho]](s),
\end{equation}
where $\delta\psi[\delta\rho]$ is the unique solution to the linearized
field equation $\nabla^2\delta\psi = -(8\pi G/c^2)\delta\rho$.
\end{definition}

\subsection{One-Way Function Equivalence}

\begin{theorem}[Geometric Authentication as One-Way Function]
\label{thm:one-way}
The source-to-fingerprint map $\Phi$ is a one-way function in the
following precise sense:
\begin{enumerate}
\item \emph{Forward evaluation:} Given $\delta\rho$, computing
$\Phi[\delta\rho]$ requires solving the Poisson equation --- achievable
in $\order{N\log N}$ time via FFT for $N$ discretization points.
\item \emph{Inverse evaluation:} Given $\Fingerprint_\Corridor$ on the
1D corridor, recovering $\delta\rho$ on $\mathbb{R}^3$ requires solving
a severely ill-posed inverse problem: the Cauchy problem for the Laplace
equation with data on a 1D curve in 3D space.
\end{enumerate}
\end{theorem}

\begin{proof}
\emph{Forward direction:} Standard. The Poisson equation
$\nabla^2\delta\psi = -S$ with $S = 8\pi G\delta\rho/c^2$ has the
explicit Green's function solution
$\delta\psi(\mathbf{r}) = \int G(\mathbf{r}-\mathbf{r}')S(\mathbf{r}')d^3r'$
with $G(\mathbf{r}) = 1/(4\pi|\mathbf{r}|)$. Discrete evaluation is
$\order{N\log N}$ via FFT-accelerated convolution.

\emph{Inverse direction:} We must show that the inverse problem ---
determining $\delta\rho(\mathbf{r})$ in $\mathbb{R}^3$ from knowledge
of $\delta\psi$ and its derivatives along a 1D curve $\Corridor$ --- is
exponentially ill-conditioned.

The Cauchy problem for the Laplace equation is ill-posed in the sense of
Hadamard~\cite{Hadamard1923}: solutions do not depend continuously on
the data. Specifically, consider the problem of determining
$\delta\psi$ in a neighborhood $\Omega$ of $\Corridor$ from Cauchy data
$(\delta\psi|_\Corridor, \partial_\nu\delta\psi|_\Corridor)$.
By the classical Hadamard instability result, the continuation from a
1D curve into 3D has condition number
\begin{equation}
\kappa_N \geq C\,e^{\pi N d/L},
\label{eq:hadamard}
\end{equation}
where $N$ is the number of Fourier modes, $d$ is the distance from
$\Corridor$ to the source region, and $L$ is the corridor length.

This means that to determine the source to precision $\epsilon$ from
corridor data known to precision $\delta$, we need
\begin{equation}
\delta < \epsilon\,e^{-\pi N d/L}.
\end{equation}
For $N = 100$ modes, $d = 10\;\text{m}$, $L = 1\;\text{km}$:
$\kappa_{100} > e^{3.14} \approx 23$. For $d \sim L$:
$\kappa_{100} > e^{314} \approx 10^{136}$, making inversion
computationally intractable.

For \emph{discrete} source configurations (point masses at unknown
positions), the inverse problem reduces to nonlinear parameter estimation.
We prove NP-hardness by reduction from the subset-sum problem.

\begin{lemma}\label{lem:np-hard}
Given $N$ possible source locations $\{\mathbf{r}_i\}$ and a target
fingerprint $\Fingerprint^*$, the problem of determining which subset
$S \subseteq \{1,\ldots,N\}$ of sources is active such that
$\|\Phi[\sum_{i\in S} m_i\delta(\mathbf{r}-\mathbf{r}_i)] -
\Fingerprint^*\| < \epsilon$ is NP-hard.
\end{lemma}

\begin{proof}[Proof of Lemma~\ref{lem:np-hard}]
Reduce from SUBSET-SUM. Given integers $\{a_1,\ldots,a_N\}$ and
target $T$, construct:
\begin{itemize}
\item $N$ sources at positions along a line perpendicular to $\Corridor$
at distances $\{d_i\}$ chosen so that
$\Phi[m_i\delta(\mathbf{r}-\mathbf{r}_i)]|_{s=L/2} = a_i/C_{\rm norm}$
for an appropriate normalization constant $C_{\rm norm}$.
\item Target fingerprint
$\Fingerprint^*(L/2) = T/C_{\rm norm}$.
\end{itemize}
Then $\|\Phi[\sum_{i\in S}m_i\delta_i] - \Fingerprint^*\| < \epsilon$
if and only if $|\sum_{i\in S}a_i - T| < \epsilon C_{\rm norm}$,
which for sufficiently small $\epsilon$ solves SUBSET-SUM.
\end{proof}
\end{proof}

\subsection{Min-Entropy of the Geometric Fingerprint}

\begin{theorem}[Fingerprint Min-Entropy]\label{thm:min-entropy}
Let the corridor fingerprint be sampled at $N_s$ points
$\{s_1,\ldots,s_{N_s}\}$ with measurement precision $\sigma_\psi$
per component. The min-entropy of the fingerprint, viewed as a
cryptographic key, is
\begin{equation}
\boxed{
\Hmin(\Fingerprint) \geq N_s \cdot 5 \cdot
\log_2\!\left(\frac{\Delta\psi_{\rm range}}{2\sigma_\psi}\right)
= 5 N_s \log_2\!\left(\frac{L}{2\ell_{\rm corr}}\cdot
\frac{\psi_{\rm max}}{\sigma_\psi}\right),
}
\label{eq:min-entropy}
\end{equation}
where $\Delta\psi_{\rm range}$ is the dynamic range of the fingerprint,
$\ell_{\rm corr}$ is the spatial correlation length of the field, and
$\psi_{\rm max}$ is the peak perturbation amplitude.
\end{theorem}

\begin{proof}
The min-entropy is defined as
$\Hmin = -\log_2\max_{\mathbf{f}} \Prob[\Fingerprint = \mathbf{f}]$.
The fingerprint at each sampling point $s_i$ takes values in
$\mathbb{R}^5$; after quantization to precision $\sigma_\psi$, each
component takes one of $N_{\rm levels} = \Delta\psi_{\rm range}/(2\sigma_\psi)$
distinguishable values. If the $N_s$ sampling points are separated by
at least $\ell_{\rm corr}$ (ensuring approximate independence), the
total number of distinguishable fingerprints is at least
$N_{\rm levels}^{5N_s}$, and the probability of any specific fingerprint
is at most $N_{\rm levels}^{-5N_s}$. Taking logarithms:
\begin{equation}
\Hmin \geq 5N_s\log_2 N_{\rm levels}
= 5N_s\log_2\!\left(\frac{\Delta\psi_{\rm range}}{2\sigma_\psi}\right).
\end{equation}
The dynamic range is $\Delta\psi_{\rm range} \sim \psi_{\rm max}$ and
the maximum number of independent samples is $N_s^{\rm max} = L/\ell_{\rm corr}$,
giving the second form.
\end{proof}

\paragraph{Numerical evaluation.}
For a corridor of length $L = 1\;\text{km}$ with correlation length
$\ell_{\rm corr} = 1\;\text{m}$ (set by the source geometry scale),
$\psi_{\rm max}/\sigma_\psi = 10^6$ (measurement SNR):
\begin{equation}
\Hmin \geq 5\times 1000\times\log_2(500\times 10^6)
= 5000\times 28.9 = 144{,}500\;\text{bits}.
\end{equation}
This is a 144.5 kbit cryptographic key equivalent --- enormously exceeding
the 256-bit keys used in AES-256 by a factor of 564.

\subsection{Adversary's Spoofing Probability}

\begin{theorem}[Exponential Spoofing Decay]\label{thm:spoof-decay}
An adversary located at displacement $d$ from the legitimate corridor
axis, attempting to spoof the fingerprint by generating a $\psi$-field
from her own source at $\mathbf{r}_E$, achieves false-acceptance
probability
\begin{equation}
\boxed{
P_{\rm spoof}(d) \leq \exp\!\left(-\frac{\beta\,N_s\,d^2}{\ell_{\rm corr}^2}\right),
}
\label{eq:spoof-decay}
\end{equation}
where the decay constant is
\begin{equation}
\beta = \frac{1}{2}\left(\frac{\psi_{\rm max}}{\sigma_\psi}\right)^2
\cdot \frac{1}{(1 + d/R_0)^6}
\label{eq:beta}
\end{equation}
with $R_0$ the emitter-corridor distance and $\sigma_\psi$ the
measurement noise.
\end{theorem}

\begin{proof}
Eve's field at corridor point $s$, generated from position
$\mathbf{r}_E = \mathbf{r}_A + \mathbf{d}$ (where $\mathbf{r}_A$ is
Alice's position and $\mathbf{d}$ is the displacement vector with
$|\mathbf{d}| = d$), differs from Alice's field by
\begin{equation}
\Delta\psi(s) \equiv \psi_E(s) - \psi_A(s)
= \frac{GM}{c^2}\left(\frac{1}{|\mathbf{r}(s) - \mathbf{r}_E|}
- \frac{1}{|\mathbf{r}(s) - \mathbf{r}_A|}\right).
\end{equation}
For $d \ll R_0$ (Eve near Alice), Taylor expansion gives
\begin{equation}
\Delta\psi(s) \approx \frac{GM}{c^2 R(s)^2}\,\hat{R}(s)\cdot\mathbf{d}
+ \order{d^2/R_0^2},
\end{equation}
where $R(s) = |\mathbf{r}(s) - \mathbf{r}_A|$ and
$\hat{R}(s) = (\mathbf{r}(s)-\mathbf{r}_A)/R(s)$.

The gradient deviation is
\begin{equation}
\Delta(\nabla\psi)(s) \approx \frac{GM}{c^2 R(s)^3}
\bigl(3\hat{R}(s)(\hat{R}(s)\cdot\mathbf{d}) - \mathbf{d}\bigr)
+ \order{d^2/R_0^3}.
\end{equation}

The 5D fingerprint deviation at each sampling point has squared norm
\begin{align}
|\Delta\Fingerprint(s_i)|^2 &= |\Delta\psi(s_i)|^2
+ |\Delta\dot\psi(s_i)|^2 + |\Delta\nabla\psi(s_i)|^2 \nonumber\\
&\geq \left(\frac{GM}{c^2 R_0^3}\right)^2 d^2
\cdot g(s_i),
\end{align}
where $g(s_i) \geq 1$ is a geometric factor from the angular variation
along the corridor.

The GLRT statistic under Eve's hypothesis is
\begin{equation}
\Lambda_E = \sum_{i=1}^{N_s}\frac{|\Delta\Fingerprint(s_i)|^2}{\sigma_\psi^2}
\geq \frac{N_s d^2}{\sigma_\psi^2}\left(\frac{GM}{c^2 R_0^3}\right)^2
\bar{g},
\end{equation}
where $\bar{g} = N_s^{-1}\sum_i g(s_i) \geq 1$.

The false-acceptance probability (probability that $\Lambda_E$ falls
below the authentication threshold $\eta$) is bounded by the
Chernoff bound on the non-central $\chi^2$ distribution:
\begin{equation}
P_{\rm spoof} \leq \exp\!\left(-\frac{\lambda_{\rm nc}}{2}
\left(1 - \sqrt{\frac{5N_s}{\lambda_{\rm nc}}}\right)^2\right),
\end{equation}
where $\lambda_{\rm nc} = \Lambda_E$. Substituting and defining
$\beta$ to absorb the geometric and SNR factors:
\begin{equation}
\beta = \frac{1}{2}\left(\frac{GM/(c^2 R_0^3)}{\sigma_\psi/\ell_{\rm corr}}\right)^2\bar{g}
= \frac{1}{2}\left(\frac{\psi_{\rm max}/R_0^2}{\sigma_\psi/\ell_{\rm corr}}\right)^2\bar{g},
\end{equation}
and noting that $\psi_{\rm max} \sim GM/(c^2 R_0)$, so
$\psi_{\rm max}/R_0^2 \sim \psi_{\rm max}/R_0^2$, we obtain
\eqref{eq:spoof-decay} with~\eqref{eq:beta}.
\end{proof}

\paragraph{Numerical evaluation.}
For $\psi_{\rm max}/\sigma_\psi = 10^6$, $R_0 = 100\;\text{m}$,
$\ell_{\rm corr} = 1\;\text{m}$, $N_s = 1000$:
\begin{itemize}
\item At $d = 1\;\text{m}$:
$P_{\rm spoof} \leq e^{-10^{12}\times 1000\times 1/1} = e^{-10^{15}}
\approx 10^{-4.3\times 10^{14}}$.
\item At $d = 1\;\text{cm}$:
$P_{\rm spoof} \leq e^{-10^{12}\times 1000\times 10^{-4}}
= e^{-10^{11}} \approx 10^{-4.3\times 10^{10}}$.
\item At $d = 1\;\mu\text{m}$:
$P_{\rm spoof} \leq e^{-10^{12}\times 1000\times 10^{-12}}
= e^{-1000} \approx 10^{-434}$.
\end{itemize}
Even at micrometer-scale displacement, the spoofing probability is
negligibly small --- beyond the square of the inverse age of the universe
in Planck times.

\subsection{Comparison with Computational Security}

\begin{proposition}[Security Equivalence]\label{prop:security-equiv}
The geometric authentication of a $\psi$-corridor with $N_s = 1000$
sampling points at SNR~$= 10^6$ provides:
\begin{itemize}
\item Min-entropy: $144{,}500$ bits $\gg 256$ bits (AES-256).
\item Spoofing probability at $d = 1\;\text{m}$:
$10^{-4.3\times 10^{14}} \ll 2^{-256} \approx 10^{-77}$ (AES-256
brute-force bound).
\item Inversion hardness: NP-hard (Lemma~\ref{lem:np-hard}),
equivalent to best-known post-quantum security.
\end{itemize}
\end{proposition}

The security margin exceeds all known cryptographic schemes by hundreds
of orders of magnitude, fundamentally because the authentication is
\emph{physical} (bound to 3D geometry) rather than \emph{computational}
(bound to algorithmic hardness assumptions).

%=============================================================================
\section{Comparison with 15 Existing Technologies}\label{sec:comparison}
%=============================================================================

\subsection{GPS/GNSS Anti-Jam Systems}

Controlled Reception Pattern Antennas (CRPAs) with 7-element arrays
provide $\sim 30\;\text{dB}$ jam resistance via null
steering~\cite{Fante2004CRPA,Hatke2015CRPA}. The fundamental limitation
is that CRPAs require $N_{\rm elements} - 1$ degrees of freedom to null
$N_{\rm elements}-1$ jammers, and performance degrades in dense jammer
environments. GPS military M-code provides $\sim 20\;\text{dB}$
additional jam resistance through spread-spectrum processing gain.

\paragraph{DFD advantage.}
$\psi$-corridors are \emph{inherently} unjammable: the information
carrier is the scalar gravitational field, not an electromagnetic
signal. No amount of EM power can disrupt $\psi$-channel reception.
This provides infinite effective jam resistance, surpassing any
antenna-based countermeasure.

\subsection{GPS Spoofing Countermeasures}

Current anti-spoofing relies on:
(i)~cryptographic authentication (GPS III L1C CHIMERA, Galileo
OSNMA)~\cite{Humphreys2012,Psiaki2016Spoofing},
(ii)~multi-antenna angle-of-arrival verification,
(iii)~inertial cross-checks, and
(iv)~power-level monitoring.
Humphreys \textit{et al.}~demonstrated that a \$3000 software-defined
radio can spoof civilian GPS, and even military P(Y) code is vulnerable
to meaconing (record-and-replay) attacks~\cite{Humphreys2012}.

\paragraph{DFD advantage.}
$\psi$-corridor authentication relies on the physical geometry of the
source mass distribution --- not on a shared cryptographic secret that
can be compromised. The spoofing impossibility theorem
(Theorem~\ref{thm:spoof-decay}) shows that the adversary must
physically replicate the mass distribution at its exact location, which
is a thermodynamic impossibility for macroscopic objects.

\subsection{Quantum Key Distribution}

BB84~\cite{Bennett1984BB84} and its variants (E91, B92, decoy-state,
SARG04, DPS, COW) derive security from the no-cloning
theorem~\cite{Wiesner1983,Scarani2009QKDreview}. Practical limitations
include:
\begin{itemize}
\item Fiber loss: $\sim 0.2\;\text{dB/km}$ limits range to $\sim 500\;\text{km}$
without quantum repeaters~\cite{Takeoka2014,Pirandola2017}.
\item Key rates: $\sim 10\;\text{Mbps}$ at short range, dropping to
$\sim 1\;\text{bps}$ at 500~km~\cite{Boaron2018}.
\item Detector attacks: Blinding~\cite{Lydersen2010Blinding},
after-gate~\cite{Wiechers2011}, and Trojan-horse
attacks~\cite{Jain2014Trojan} exploit implementation imperfections.
\item Device-independent QKD addresses detector vulnerabilities but at
the cost of much lower key rates ($\sim 10\;\text{bps}$).
\end{itemize}

\paragraph{DFD advantage.}
$\psi$-corridor security is:
(i)~distance-independent (no fiber loss analogue),
(ii)~immune to detector-side attacks (the fingerprint is a
classical macroscopic observable),
(iii)~provides vastly higher key equivalent
($\Hmin \sim 10^5$ bits vs.\ typical QKD session key of $\sim 10^6$
bits, but at unlimited range).
The security basis differs: QKD uses quantum no-cloning;
$\psi$-authentication uses elliptic PDE inversion hardness.

\subsection{Neutrino Communications}

The Fermilab NuMI experiment~\cite{Stancil2012Neutrino} demonstrated
digital communication through 240~m of rock using the NuMI neutrino
beam and the MINERvA detector, achieving a data rate of
$\sim 0.1\;\text{bits/s}$ at BER $\sim 1\%$. The ANTARES
collaboration proposed submarine neutrino communication
links~\cite{Ageron2011ANTARES}. Fundamental limitations: neutrinos
require particle accelerators (multi-MW beam power) and kiloton-scale
detectors, making practical deployment impossible with current technology.

\paragraph{DFD advantage.}
$\psi$-corridors penetrate all media (like neutrinos) but with
dramatically lower infrastructure requirements: no particle accelerator,
no kiloton detector. The quasi-static nature of the $\psi$-field means
the ``signal'' does not require particle production and detection.

\subsection{Underwater Acoustic Communications}

State-of-the-art underwater acoustic modems achieve
$\sim 50\;\text{kbps}$ at 1--5~km range, with severe multipath
($\sim 10\;\text{ms}$ spread), Doppler ($v/c_{\rm water} \sim 10^{-3}$),
and frequency-dependent absorption ($\alpha \sim f^2/(1+f^2/f_0^2)$
dB/km)~\cite{Stojanovic2009Underwater,Chitre2008Underwater}. The
capacity of the underwater acoustic channel is fundamentally limited by
the Wenz curves of ambient ocean noise and the narrow usable bandwidth
($\sim 10$--$50\;\text{kHz}$)~\cite{Ainslie2010Principles}.
Recent work on MIMO acoustic systems and OFDM modulation pushes rates
toward 100~kbps~\cite{Li2018MIMO_UW}, but kilometer-range high-rate
communication remains elusive.

\paragraph{DFD advantage.}
$\psi$-corridors propagate through seawater without acoustic or EM
attenuation. There is no multipath (the quasi-static field follows the
source geometry, not reflected paths), no Doppler (the field is
quasi-static), and no frequency-dependent absorption. The achievable
bandwidth is limited only by the $\psi$-emitter modulation rate, not by
the medium.

\subsection{Through-Wall Communications and Radar}

Ultra-wideband (UWB) through-wall radar operates at 1--10~GHz with
typical wall penetration loss of 5--30~dB depending on material
(drywall: 5~dB; concrete: 15--25~dB; reinforced concrete: 25--40~dB)
\cite{Amin2011ThroughWall,Ahmad2012ThroughWall}. The round-trip path
loss through two walls plus free-space propagation severely limits
range. Through-wall communication via UWB achieves $\sim 1\;\text{Mbps}$
through standard residential construction but fails through
concrete/steel structures.

\paragraph{DFD advantage.}
$\psi$-fields penetrate all materials without attenuation
(the scalar field couples to mass-energy, not EM properties). There is
no wall penetration loss, enabling communication through reinforced
concrete, underground, and into shielded enclosures (Faraday cages).

\subsection{Military Anti-Jam Communication Systems}

The MILSATCOM anti-jam architecture includes:
(i)~Advanced EHF (AEHF) with $\sim 60\;\text{dB}$ jam resistance via
frequency hopping and nulling antennas~\cite{Warner2019AEHF},
(ii)~Link-16/JTIDS with pseudo-random frequency-hopping TDMA,
(iii)~Laser SATCOM (LCRD) providing inherent jam resistance via narrow
beamwidth, and
(iv)~Protected Tactical Waveform (PTW) for tactical
networks. These systems achieve anti-jam margins of 30--60~dB but at
enormous cost (\$1--10B per satellite) and with the fundamental
limitation that they remain electromagnetic, hence vulnerable to
sufficiently powerful jammers.

\paragraph{DFD advantage.}
$\psi$-corridors provide \emph{absolute} jam immunity at any jammer
power level. The non-electromagnetic nature of the carrier eliminates
the arms race between jammer power and anti-jam processing gain.

\subsection{Quantum Navigation Proposals}

Atom interferometry for inertial navigation achieves acceleration
sensitivity $\sim 10\;\text{nGal}$ and rotation sensitivity
$\sim 10^{-9}\;\text{rad/s}$~\cite{Geiger2020AtomInterferometry,
Narducci2022QuantumNav}. Cold-atom gravimeters have demonstrated
$\sim 10^{-8}g$ sensitivity in field conditions. Fundamental
limitations: atom interferometers measure \emph{inertial} quantities
(acceleration, rotation) that must be integrated, accumulating drift.
Absolute position requires an external reference (gravity map matching).

\paragraph{DFD advantage.}
$\psi$-field positioning provides \emph{absolute} position directly
from the field geometry, without integration drift. The $\psi$-field is
the gravitational field itself, so no gravity map matching is needed ---
the ``map'' is the real-time field measurement.

\subsection{Inertial Navigation Drift}

Navigation-grade ring laser gyros achieve $\sim 0.01\;\text{deg/hr}$
drift, giving position errors of $\sim 1\;\text{nm/hr}$ (nautical
mile per hour)~\cite{Chatfield1997,Titterton2004INS}. Strategic-grade
INS (submarine, ICBM) achieves $\sim 0.001\;\text{deg/hr}$ with
position drift $\sim 0.1\;\text{nm/hr}$. MEMS INS drifts at
$\sim 10\;\text{nm/hr}$. All INS drift without bound and require
external aiding.

\paragraph{DFD advantage.}
$\psi$-positioning has zero drift: each measurement is an independent
absolute position fix. The system can provide continuous drift-free
aiding to any INS, analogous to GPS but without EM vulnerability.

\subsection{Summary Comparison Table}

\begin{table*}[t]
\centering
\caption{Comprehensive comparison of $\psi$-corridor communications
with 15 existing technologies. ``Distance'' refers to practical operating
range. ``Auth'' indicates native authentication capability.
``Jam'' indicates jam resistance.}
\label{tab:comprehensive}
\begin{ruledtabular}
\begin{tabular}{lcccccl}
Technology & Rate & Distance & Penetrate & Auth & Jam & Key Limitation \\
\hline
5G mmWave & 10 Gbps & 200 m & No & No & No & EM jamming, LOS \\
FSO (1550 nm) & 100 Gbps & 10 km & No & No & No & Weather, LOS \\
Fiber optical & 100 Tbps & 10000 km & N/A & No & Hard & Physical tap \\
QKD (BB84) & 10 kbps & 500 km & No & Yes & No & Distance, rate \\
Neutrino & 0.1 bps & Earth & Yes & No & Yes & MW accelerator \\
UW acoustic & 50 kbps & 5 km & Water & No & Mod & Multipath, BW \\
UW EM & 1 kbps & 10 m & Water & No & No & Skin depth \\
UWB through-wall & 1 Mbps & 30 m & Partial & No & Low & Material loss \\
GPS L1C & 50 bps & Global & No & Weak & No & Jamming, spoofing \\
GPS M-code & 50 bps & Global & No & Yes & 20 dB & Cost, access \\
AEHF MILSATCOM & 8 Mbps & Global & No & Yes & 60 dB & \$10B/sat \\
Link-16 JTIDS & 1 Mbps & 500 km & No & Yes & 30 dB & LOS, cost \\
Atom interf.\ nav & --- & N/A & N/A & N/A & Yes & Drift, SWaP \\
RLG INS & --- & N/A & N/A & N/A & Yes & Drift 1 nm/hr \\
Gravity gradiometry & --- & Local & Yes & N/A & Yes & Sensitivity \\
\hline
$\psi$-corridor & $5B\log_2(1\!+\!\SNR/5)$ & Any & Yes & Yes & $\infty$ & $\psi$-pump power \\
\end{tabular}
\end{ruledtabular}
\end{table*}

%=============================================================================
\section{Combined Communications + Navigation Demonstrator}
\label{sec:demonstrator}
%=============================================================================

\subsection{Unified Infrastructure Concept}

A single $\psi$-corridor simultaneously provides three services:
\begin{enumerate}
\item \textbf{Data communication:} Modulated $\delta\psi(t)$ carries
information bits.
\item \textbf{Positioning:} The corridor geometry (gradient map) enables
receiver position determination.
\item \textbf{Authentication:} The corridor fingerprint provides identity
verification.
\end{enumerate}
These services share the same physical infrastructure ($\psi$-emitters
and corridors) but compete for the finite $\psi$-modulation bandwidth.

\subsection{Resource Allocation Framework}

Let $B$ be the total $\psi$-modulation bandwidth and $P$ the total
$\psi$-pump power. The three services have bandwidth allocations
$B_c$, $B_n$, $B_a$ with $B_c + B_n + B_a \leq B$, and power
allocations $P_c$, $P_n$, $P_a$ with $P_c + P_n + P_a \leq P$.

\paragraph{Communication capacity:}
\begin{equation}
C = 5\,B_c\log_2\!\left(1 + \frac{P_c}{5\,N_0\,B_c}\right).
\end{equation}

\paragraph{Navigation precision:}
The Cram\'er-Rao bound on position error scales inversely with
the SNR devoted to navigation:
\begin{equation}
\sigma_{\rm pos}^2 \geq \frac{N_0\,B_n}{\text{GDOP}^{-2}\cdot P_n\cdot B_n\cdot T_n}
= \frac{N_0}{\text{GDOP}^{-2}\cdot P_n \cdot T_n},
\end{equation}
where $T_n$ is the navigation integration time and GDOP depends on
emitter geometry.

\paragraph{Authentication strength:}
The false-acceptance probability scales as
\begin{equation}
P_{\rm FA} \leq \exp\!\left(-\frac{P_a\,T_a}{N_0\,B_a}\cdot N_s\right)
\end{equation}
where $T_a$ is the authentication integration time.

\subsection{Convex Optimization of Resource Allocation}

\begin{theorem}[Optimal Resource Allocation]
\label{thm:resource}
The problem of maximizing communication rate $C$ subject to minimum
positioning precision $\sigma_{\rm pos} \leq \sigma_{\rm max}$ and
minimum authentication security $P_{\rm FA} \leq P_{\rm FA}^{\rm max}$
is a convex optimization problem:
\begin{align}
\max_{B_c,B_n,B_a,P_c,P_n,P_a} &\quad 5B_c\log_2\!\left(1 + \frac{P_c}{5N_0 B_c}\right) \label{eq:opt-obj}\\
\text{s.t.}\quad & \frac{N_0}{P_n T_n}\cdot\text{GDOP}^2 \leq \sigma_{\rm max}^2, \label{eq:opt-nav}\\
& \frac{P_a T_a N_s}{N_0 B_a} \geq \ln(1/P_{\rm FA}^{\rm max}), \label{eq:opt-auth}\\
& B_c + B_n + B_a \leq B, \label{eq:opt-bw}\\
& P_c + P_n + P_a \leq P, \label{eq:opt-power}\\
& B_c, B_n, B_a, P_c, P_n, P_a \geq 0. \label{eq:opt-nonneg}
\end{align}
\end{theorem}

\begin{proof}
The objective~\eqref{eq:opt-obj} is concave in $(B_c, P_c)$ (composition
of concave $\log$ with affine argument). The constraint~\eqref{eq:opt-nav}
is linear in $P_n$. The constraint~\eqref{eq:opt-auth} is the ratio
$P_a/(N_0 B_a) \geq \text{const}$; after taking $B_a = B_a^{\rm fixed}$
(which can be optimized separately), this is linear in $P_a$. All
remaining constraints are linear. The problem is therefore convex and
can be solved by KKT conditions or interior-point methods.
\end{proof}

\paragraph{KKT solution.}
The Lagrangian is
$\mathcal{L} = C - \lambda_1(g_1 - \sigma_{\rm max}^2) - \lambda_2(h_2 - g_2) - \mu_B(B_c+B_n+B_a-B) - \mu_P(P_c+P_n+P_a-P)$.
Setting $\partial\mathcal{L}/\partial P_c = 0$:
\begin{equation}
\frac{5B_c}{(P_c + 5N_0 B_c)\ln 2} = \mu_P,
\end{equation}
giving $P_c = 5B_c/(\mu_P\ln 2) - 5N_0 B_c$.
The navigation constraint gives $P_n = N_0\,\text{GDOP}^2/(\sigma_{\rm max}^2 T_n)$.
The authentication constraint gives
$P_a = N_0 B_a\ln(1/P_{\rm FA}^{\rm max})/(T_a N_s)$.
The remaining power goes to communication:
$P_c = P - P_n - P_a$.

\subsection{Demonstrator Design Parameters}

\begin{table}[t]
\caption{Combined comms+nav demonstrator parameters.}
\label{tab:demonstrator}
\begin{ruledtabular}
\begin{tabular}{lcc}
Parameter & Symbol & Value \\
\hline
Total bandwidth & $B$ & 10 kHz \\
Comm bandwidth & $B_c$ & 8 kHz \\
Nav bandwidth & $B_n$ & 1 kHz \\
Auth bandwidth & $B_a$ & 1 kHz \\
Total $\psi$-pump power & $P$ & $10^{-38}\;\text{s}^{-1}$ \\
Comm power fraction & $P_c/P$ & 0.90 \\
Nav power fraction & $P_n/P$ & 0.05 \\
Auth power fraction & $P_a/P$ & 0.05 \\
Detector volume & $V_d$ & $(1\;\text{m})^3$ \\
Operating range & $R$ & 100 m \\
Number of emitters & $N_E$ & 4 \\
Auth sampling points & $N_s$ & 100 \\
Nav integration time & $T_n$ & 1 s \\
Auth integration time & $T_a$ & 10 s \\
\end{tabular}
\end{ruledtabular}
\end{table}

\paragraph{Predicted performance.}
With the parameters in Table~\ref{tab:demonstrator}:
\begin{itemize}
\item Communication rate:
$C \approx 5\times 8\times 10^3\times\log_2(1 + 10^{29}/5)
\approx 40\times 10^3\times 96 \approx 3.8\;\text{Mbps}$ (theoretical maximum).
\item Nav precision:
$\sigma_{\rm pos} \approx \text{GDOP}\times\sqrt{N_0/(P_n T_n)}$;
for GDOP$\sim 2$ and $P_n T_n/N_0 \sim 10^{27}$:
$\sigma_{\rm pos} \sim 10^{-13}\;\text{m}$ (sub-femtometer theoretical
limit --- practically limited by emitter stability).
\item Auth security:
$P_{\rm FA} \leq e^{-P_a T_a N_s/(N_0 B_a)}
\leq e^{-10^{27}\times 100/10^3} = e^{-10^{26}}$ ---
cryptographically impregnable.
\end{itemize}

\subsection{Practical Demonstrator Architecture}

The demonstrator consists of:
\begin{enumerate}
\item \textbf{Four $\psi$-emitters} at the corners of a
$100\;\text{m}\times 100\;\text{m}$ square, each comprising a
dense rotating mass (tungsten rotor array) modulated at distinct
frequencies $\{f_1, f_2, f_3, f_4\} = \{100, 200, 300, 400\}\;\text{Hz}$.
\item \textbf{Central receiver} with a high-Q cavity $\psi$-detector
(optical clock-based or atom interferometer) and a multi-axis
gradiometer array.
\item \textbf{Signal processing} unit performing:
\begin{itemize}
\item Frequency-domain separation of emitter signals,
\item 5D constellation demodulation for data recovery,
\item Gradient-based trilateration for position fix,
\item Fingerprint correlation for authentication.
\end{itemize}
\end{enumerate}

The key technological challenge is achieving sufficient $\psi$-pump
amplitude. The EM--$\psi$ coupling parameter
$|\lambda - 1| \lesssim 3\times 10^{-5}$ sets the conversion efficiency
from electromagnetic to $\psi$-field energy. Intentional resonant
enhancement in high-Q cavities can increase the effective coupling.

%=============================================================================
\section{Information-Theoretic Bounds on Combined System}
\label{sec:bounds}
%=============================================================================

\subsection{Joint Communication-Navigation-Authentication Rate Region}

\begin{definition}[Rate Region]
The achievable rate region $\mathcal{R}$ of the combined system is the
set of triples $(C, \sigma_{\rm pos}^{-2}, -\log P_{\rm FA})$
simultaneously achievable subject to the total resource constraints.
\end{definition}

\begin{theorem}[Rate Region Characterization]
The rate region boundary is parameterized by
$(\alpha_c, \alpha_n, \alpha_a)$ with $\alpha_c + \alpha_n + \alpha_a = 1$,
$\alpha_i \geq 0$:
\begin{align}
C(\alpha_c) &= 5\alpha_c B\log_2\!\left(1 + \frac{\alpha_c P}{5N_0\alpha_c B}\right)
= 5\alpha_c B\log_2\!\left(1 + \frac{P}{5N_0 B}\right), \\
\sigma_{\rm pos}^{-2}(\alpha_n) &= \frac{\alpha_n P\,T_n}{N_0}\,\text{GDOP}^{-2}, \\
-\log P_{\rm FA}(\alpha_a) &= \frac{\alpha_a P\,T_a\,N_s}{N_0\,\alpha_a B}
= \frac{P\,T_a\,N_s}{N_0\,B}.
\end{align}
\end{theorem}

Remarkably, the authentication security $-\log P_{\rm FA}$ is
\emph{independent} of the bandwidth allocation $\alpha_a$ (the ratio
$P_a/B_a = \alpha_a P/(\alpha_a B) = P/B$ is constant). This means
authentication can be provided ``for free'' at any desired security level
simply by increasing the integration time $T_a$, without reducing
communication capacity or navigation precision.

%=============================================================================
\section{Conclusions}\label{sec:conclusions}
%=============================================================================

We have derived the complete Shannon capacity of $\psi$-field
communication channels from first principles, starting with the quantum
vacuum noise floor
$N_0 = 4\hbar G/(c^3 V_d) \approx 10^{-69}\;\text{Hz}^{-1}$. The 5D
composite channel achieves a capacity gain approaching $5\times$ over
the scalar channel at high SNR, with explicit sphere-packing analysis
confirming the geometric advantage.

The central security result is the proof that geometry-locked
authentication via corridor fingerprinting is computationally equivalent
to inverting an elliptic PDE --- an NP-hard problem for discrete source
configurations (Theorem~\ref{thm:one-way}). The min-entropy of the
geometric fingerprint reaches $\sim 144{,}500$ bits for a 1~km corridor,
and the adversary's spoofing probability decays as
$\exp(-\beta N_s d^2/\ell_{\rm corr}^2)$ with displacement $d$
(Theorem~\ref{thm:spoof-decay}).

Comparison with 15 existing technologies confirms that $\psi$-corridors
uniquely combine medium penetration, native authentication, and
absolute jam immunity --- capabilities not simultaneously available in
any other communication modality.

The combined communications+navigation+authentication demonstrator
shows that a single $\psi$-corridor infrastructure can deliver all three
services with convex-optimal resource allocation, and that authentication
is achievable with negligible resource overhead.

%=============================================================================
\begin{thebibliography}{99}

\bibitem{Alcock2025DFD}
G.~T. Alcock, ``Density Field Dynamics: A Complete Unified Theory,''
v3.2 (2025).

\bibitem{Alcock2025PsiComms}
G.~T. Alcock, ``$\psi$-Field Communications: Geometry-Locked,
Spoof-Proof Information Transfer via Scalar Refractive Corridors,''
DFD Research Programme (2025).

\bibitem{Shannon1948}
C.~E. Shannon, ``A Mathematical Theory of Communication,''
Bell Syst.\ Tech.\ J.\ \textbf{27}, 379--423, 623--656 (1948).

\bibitem{CoverThomas2006}
T.~M. Cover and J.~A. Thomas, \textit{Elements of Information Theory},
2nd ed.\ (Wiley-Interscience, 2006).

\bibitem{Hadamard1923}
J.~Hadamard, \textit{Lectures on Cauchy's Problem in Linear Partial
Differential Equations} (Yale University Press, 1923).

\bibitem{Humphreys2012}
T.~E. Humphreys, B.~M. Ledvina, M.~L. Psiaki, B.~W. O'Hanlon, and
P.~M. Kintner, ``Assessing the Spoofing Threat: Development of a
Portable GPS Civilian Spoofer,'' in \textit{Proc.\ ION GNSS}
(2008), pp.\ 2314--2325.

\bibitem{Psiaki2016Spoofing}
M.~L. Psiaki and T.~E. Humphreys, ``GNSS Spoofing and Detection,''
Proc.\ IEEE \textbf{104}, 1258--1270 (2016).

\bibitem{Bennett1984BB84}
C.~H. Bennett and G.~Brassard, ``Quantum Cryptography: Public Key
Distribution and Coin Tossing,'' in \textit{Proc.\ IEEE Int.\ Conf.\
Comput.\ Syst.\ Signal Process.}\ (1984), pp.\ 175--179.

\bibitem{Wiesner1983}
S.~Wiesner, ``Conjugate Coding,'' SIGACT News \textbf{15}, 78--88 (1983).

\bibitem{Scarani2009QKDreview}
V.~Scarani \textit{et al.}, ``The Security of Practical Quantum Key
Distribution,'' Rev.\ Mod.\ Phys.\ \textbf{81}, 1301--1350 (2009).

\bibitem{Takeoka2014}
M.~Takeoka, S.~Guha, and M.~M. Wilde, ``Fundamental Rate-Loss
Tradeoff for Optical Quantum Key Distribution,''
Nat.\ Commun.\ \textbf{5}, 5235 (2014).

\bibitem{Pirandola2017}
S.~Pirandola \textit{et al.}, ``Fundamental Limits of Repeaterless
Quantum Communications,'' Nat.\ Commun.\ \textbf{8}, 15043 (2017).

\bibitem{Boaron2018}
A.~Boaron \textit{et al.}, ``Secure Quantum Key Distribution over 421 km
of Optical Fiber,'' Phys.\ Rev.\ Lett.\ \textbf{121}, 190502 (2018).

\bibitem{Lydersen2010Blinding}
L.~Lydersen \textit{et al.}, ``Hacking Commercial Quantum Cryptography
Systems by Tailored Bright Illumination,'' Nat.\ Photonics \textbf{4},
686--689 (2010).

\bibitem{Wiechers2011}
C.~Wiechers \textit{et al.}, ``After-Gate Attack on a Quantum
Cryptosystem,'' New J.\ Phys.\ \textbf{13}, 013043 (2011).

\bibitem{Jain2014Trojan}
N.~Jain \textit{et al.}, ``Trojan-Horse Attacks Threaten the Security
of Practical Quantum Cryptography,'' New J.\ Phys.\ \textbf{16},
123030 (2014).

\bibitem{Stancil2012Neutrino}
D.~D. Stancil \textit{et al.}, ``Demonstration of Communication Using
Neutrinos,'' Mod.\ Phys.\ Lett.\ A \textbf{27}, 1250077 (2012).

\bibitem{Ageron2011ANTARES}
M.~Ageron \textit{et al.} (ANTARES Collaboration), ``ANTARES: The First
Undersea Neutrino Telescope,'' Nucl.\ Instrum.\ Methods Phys.\ Res.\
A \textbf{656}, 11--38 (2011).

\bibitem{Stojanovic2009Underwater}
M.~Stojanovic and J.~Preisig, ``Underwater Acoustic Communication
Channels: Propagation Models and Statistical Characterization,''
IEEE Commun.\ Mag.\ \textbf{47}, 84--89 (2009).

\bibitem{Chitre2008Underwater}
M.~Chitre, S.~Shahabudeen, and M.~Stojanovic, ``Underwater Acoustic
Communications and Networking: Recent Advances and Future Challenges,''
Mar.\ Technol.\ Soc.\ J.\ \textbf{42}, 103--116 (2008).

\bibitem{Ainslie2010Principles}
M.~A. Ainslie, \textit{Principles of Sonar Performance Modelling}
(Springer, 2010).

\bibitem{Li2018MIMO_UW}
B.~Li \textit{et al.}, ``MIMO-OFDM for High-Rate Underwater Acoustic
Communications,'' IEEE J.\ Ocean.\ Eng.\ \textbf{34}, 634--644 (2009).

\bibitem{Amin2011ThroughWall}
M.~G. Amin (Ed.), \textit{Through-the-Wall Radar Imaging} (CRC Press, 2011).

\bibitem{Ahmad2012ThroughWall}
F.~Ahmad and M.~G. Amin, ``Through-the-Wall Human Motion Indication
Using Sparsity-Driven Change Detection,'' IEEE Trans.\ Geosci.\
Remote Sens.\ \textbf{51}, 881--890 (2013).

\bibitem{Fante2004CRPA}
R.~L. Fante and J.~J. Vaccaro, ``Wideband Cancellation of
Interference in a GPS Receive Array,'' IEEE Trans.\ Aerosp.\ Electron.\
Syst.\ \textbf{36}, 549--564 (2000).

\bibitem{Hatke2015CRPA}
G.~F. Hatke \textit{et al.}, ``GPS Anti-Jam Technology,''
in \textit{Proc.\ IEEE Radar Conf.}\ (2015).

\bibitem{Warner2019AEHF}
J.~S. Warner and R.~G. Johnston, ``GPS Spoofing Countermeasures,''
Homeland Secur.\ J.\ (2019).

\bibitem{Geiger2020AtomInterferometry}
R.~Geiger \textit{et al.}, ``High-Accuracy Inertial Measurements with
Cold-Atom Sensors,'' AVS Quantum Sci.\ \textbf{2}, 024702 (2020).

\bibitem{Narducci2022QuantumNav}
F.~A. Narducci \textit{et al.}, ``Advances in Atom Interferometry for
Navigation,'' in \textit{Proc.\ IEEE/ION PLANS}\ (2022).

\bibitem{Chatfield1997}
A.~B. Chatfield, \textit{Fundamentals of High Accuracy Inertial
Navigation} (AIAA, 1997).

\bibitem{Titterton2004INS}
D.~Titterton and J.~Weston, \textit{Strapdown Inertial Navigation
Technology}, 2nd ed.\ (IET, 2004).

\end{thebibliography}

\end{document}
